\subsection{From the Plane to Three-Dimensional Vectors}

Adding a third coordinate does not change the basic meaning of a vector. A
vector still records a displacement, but now it records three independent
coordinate changes:
\[
\mathbf v=(v_1,v_2,v_3).
\]
Addition and scalar multiplication remain componentwise, because successive
changes in each coordinate direction still combine independently.

Applying the Pythagorean theorem first to the first two coordinate changes and
then to the third gives
\[
\boxed{
\|\mathbf v\|
=
\sqrt{v_1^2+v_2^2+v_3^2}
}.
\]
The dot product extends in exactly the same way:
\[
\boxed{
\mathbf u\cdot\mathbf v
=
u_1v_1+u_2v_2+u_3v_3
}.
\]
Its geometric role is unchanged:
\[
\mathbf u\cdot\mathbf v
=
\|\mathbf u\|\,\|\mathbf v\|\cos\theta.
\]
Consequently,
\[
\boxed{
\mathbf u\perp\mathbf v
\iff
\mathbf u\cdot\mathbf v=0
}.
\]
Thus the dot product detects perpendicularity in both two and three dimensions.

\subsubsection*{A second product available in three dimensions}

In three dimensions there is another natural question. Given two vectors
\(\mathbf u\) and \(\mathbf v\), can we construct a new vector perpendicular to
both of them?

\begingroup
\setlength{\fboxsep}{9pt}

\begin{center}
\fcolorbox{blue!65!black}{blue!4}{%
\begin{minipage}{0.90\textwidth}
\textcolor{blue!65!black}{\textbf{Step 1: translate the geometric requirement into equations.}}

Let
\[
\mathbf u=(u_1,u_2,u_3),
\qquad
\mathbf v=(v_1,v_2,v_3),
\qquad
\mathbf n=(n_1,n_2,n_3).
\]
We require \(\mathbf n\) to be perpendicular to both inputs:
\[
\mathbf n\cdot\mathbf u=0,
\qquad
\mathbf n\cdot\mathbf v=0.
\]
Therefore
\[
\begin{aligned}
u_1n_1+u_2n_2+u_3n_3&=0,\\
v_1n_1+v_2n_2+v_3n_3&=0.
\end{aligned}
\]
\end{minipage}}
\end{center}

\begin{center}
\fcolorbox{violet!65!black}{violet!4}{%
\begin{minipage}{0.90\textwidth}
\textcolor{violet!65!black}{\textbf{Step 2: choose one component so that the remaining system is easy.}}

Set
\[
\textcolor{violet!75!black}{\Delta=u_1v_2-u_2v_1}.
\]
Assume first that \(\Delta\neq0\), and choose
\[
\textcolor{violet!75!black}{n_3=\Delta}.
\]
The two equations reduce to the ordinary \(2\times2\) system
\[
\begin{aligned}
u_1n_1+u_2n_2&=-u_3\Delta,\\
v_1n_1+v_2n_2&=-v_3\Delta.
\end{aligned}
\]
\end{minipage}}
\end{center}

\begin{center}
\fcolorbox{orange!80!black}{orange!5}{%
\begin{minipage}{0.90\textwidth}
\textcolor{orange!80!black}{\textbf{Step 3: eliminate \(n_2\) and find \(n_1\).}}

Multiply the first equation by \(v_2\), the second by \(u_2\), and subtract:
\[
\begin{aligned}
&v_2(u_1n_1+u_2n_2)
-u_2(v_1n_1+v_2n_2)\\
&\qquad=-u_3v_2\Delta+u_2v_3\Delta.
\end{aligned}
\]
The \(n_2\)-terms cancel, so
\[
\Delta n_1=\Delta(u_2v_3-u_3v_2).
\]
Since \(\Delta\neq0\),
\[
\boxed{\textcolor{orange!80!black}{n_1=u_2v_3-u_3v_2}}.
\]
\end{minipage}}
\end{center}

\begin{center}
\fcolorbox{orange!80!black}{orange!5}{%
\begin{minipage}{0.90\textwidth}
\textcolor{orange!80!black}{\textbf{Step 4: eliminate \(n_1\) and find \(n_2\).}}

Multiply the first equation by \(v_1\), the second by \(u_1\), and subtract:
\[
\begin{aligned}
&v_1(u_1n_1+u_2n_2)
-u_1(v_1n_1+v_2n_2)\\
&\qquad=-u_3v_1\Delta+u_1v_3\Delta.
\end{aligned}
\]
The \(n_1\)-terms cancel, giving
\[
-\Delta n_2=\Delta(u_1v_3-u_3v_1).
\]
Hence
\[
\boxed{\textcolor{orange!80!black}{n_2=u_3v_1-u_1v_3}}.
\]
\end{minipage}}
\end{center}

\begin{center}
\fcolorbox{green!55!black}{green!4}{%
\begin{minipage}{0.90\textwidth}
\textcolor{green!45!black}{\textbf{Step 5: assemble the three components.}}

Together with
\[
\textcolor{violet!75!black}{n_3=u_1v_2-u_2v_1},
\]
we obtain
\[
\boxed{
\mathbf n
=
\bigl(
\textcolor{orange!80!black}{u_2v_3-u_3v_2},
\;\textcolor{orange!80!black}{u_3v_1-u_1v_3},
\;\textcolor{violet!75!black}{u_1v_2-u_2v_1}
\bigr)
}.
\]
If the displayed \(\Delta\) is zero but the vectors are not parallel, one of the
other two coordinate pairs is nonzero. Repeating the same elimination with that
pair gives the same cyclic formula.
\end{minipage}}
\end{center}

\begin{center}
\fcolorbox{gray!70!black}{gray!5}{%
\begin{minipage}{0.90\textwidth}
\textcolor{gray!70!black}{\textbf{Step 6: verify that the constructed vector really works.}}

For the first input vector, the terms cancel in three visible pairs:
\[
\begin{aligned}
\mathbf n\cdot\mathbf u
&=\textcolor{blue!65!black}{(u_1u_2v_3-u_1u_2v_3)}\\
&\quad+\textcolor{orange!80!black}{(-u_1u_3v_2+u_1u_3v_2)}\\
&\quad+\textcolor{green!45!black}{(u_2u_3v_1-u_2u_3v_1)}\\
&=0.
\end{aligned}
\]
Likewise,
\[
\begin{aligned}
\mathbf n\cdot\mathbf v
&=\textcolor{blue!65!black}{(u_2v_1v_3-u_2v_1v_3)}\\
&\quad+\textcolor{orange!80!black}{(-u_3v_1v_2+u_3v_1v_2)}\\
&\quad+\textcolor{green!45!black}{(-u_1v_2v_3+u_1v_2v_3)}\\
&=0.
\end{aligned}
\]
Thus \(\mathbf n\perp\mathbf u\) and \(\mathbf n\perp\mathbf v\).
\end{minipage}}
\end{center}

\endgroup

The vector has therefore been derived from the two perpendicularity equations.
We now give it a name:
\[
\boxed{
\mathbf u\times\mathbf v
=
\bigl(
u_2v_3-u_3v_2,
\;u_3v_1-u_1v_3,
\;u_1v_2-u_2v_1
\bigr)
}.
\]
Consequently,
\[
\boxed{
(\mathbf u\times\mathbf v)\cdot\mathbf u=0,
\qquad
(\mathbf u\times\mathbf v)\cdot\mathbf v=0
}.
\]

The statement is stronger. Every linear combination
\[
a\mathbf u+b\mathbf v
\]
is also perpendicular to \(\mathbf u\times\mathbf v\), because
\[
\begin{aligned}
(\mathbf u\times\mathbf v)\cdot
(a\mathbf u+b\mathbf v)
&=
a(\mathbf u\times\mathbf v)\cdot\mathbf u
+b(\mathbf u\times\mathbf v)\cdot\mathbf v\\
&=a\cdot0+b\cdot0\\
&=0.
\end{aligned}
\]
We do not yet turn this family of linear combinations into a theory of planes.
That geometric construction belongs to the later chapter on lines and planes.
For now, the important fact is simply that one cross product produces a vector
perpendicular to every combination of its two inputs.

\begin{center}
\begin{tikzpicture}[scale=0.9,>=stealth]
    \coordinate (O) at (0,0);
    \coordinate (U) at (4.4,0.8);
    \coordinate (V) at (1.5,2.7);
    \coordinate (S) at (5.9,3.5);
    \coordinate (N) at (2.2,4.8);

    \draw[blue!70!black,line width=1.1pt,->] (O) -- (U)
        node[midway,below] {$\mathbf u$};
    \draw[orange!85!black,line width=1.1pt,->] (O) -- (V)
        node[midway,left] {$\mathbf v$};
    \draw[helper,dashed] (U) -- (S);
    \draw[helper,dashed] (V) -- (S);
    \draw[very thick,->] (O) -- (N)
        node[midway,left] {$\mathbf u\times\mathbf v$};

    \node[align=left,anchor=west] at (4.8,4.5)
        {perpendicular to $\mathbf u$\\
         perpendicular to $\mathbf v$\\
         perpendicular to $a\mathbf u+b\mathbf v$};
\end{tikzpicture}
\end{center}

\subsubsection*{The cross product detects parallelism}

If \(\mathbf u\) and \(\mathbf v\) are parallel, then one is a scalar multiple
of the other. Let \(\mathbf v=\lambda\mathbf u\). Substitution into the
coordinate formula makes every component of the cross product vanish, so
\[
\mathbf u\times\mathbf v=\mathbf0.
\]
Conversely, if two nonzero vectors have zero cross product, then their
directions are parallel. Hence
\[
\boxed{
\mathbf u\parallel\mathbf v
\iff
\mathbf u\times\mathbf v=\mathbf0
}
\]
for nonzero vectors.

This gives two complementary geometric tests:
\[
\boxed{
\begin{aligned}
\mathbf u\cdot\mathbf v=0
&\quad\Longleftrightarrow\quad
\mathbf u\perp\mathbf v,\\
\mathbf u\times\mathbf v=\mathbf0
&\quad\Longleftrightarrow\quad
\mathbf u\parallel\mathbf v
\quad(\mathbf u,\mathbf v\neq\mathbf0).
\end{aligned}
}
\]

\subsubsection*{Order and orientation}

There are two opposite directions perpendicular to both nonparallel vectors. To
choose one consistently, we use the right-hand rule: curl the fingers of the
right hand from \(\mathbf u\) toward \(\mathbf v\); the thumb indicates the
direction of \(\mathbf u\times\mathbf v\).

Reversing the order reverses the result:
\[
\boxed{
\mathbf v\times\mathbf u
=-(\mathbf u\times\mathbf v)
}.
\]
Thus the cross product is not commutative.

Its length satisfies
\[
\|\mathbf u\times\mathbf v\|
=
\|\mathbf u\|\,\|\mathbf v\|\sin\theta.
\]
This formula explains the parallelism test: for parallel vectors,
\(\theta=0\) or \(\theta=\pi\), so the sine and therefore the cross product
vanish. The later geometry chapters will use the same quantity when areas,
lines, and planes are constructed.

\begin{summarybox}
The two vector products answer different geometric questions:
\[
\begin{array}{c|c|c}
\text{operation}&\text{output}&\text{main test}\\
\hline
\mathbf u\cdot\mathbf v&\text{number}&
\mathbf u\perp\mathbf v\iff\mathbf u\cdot\mathbf v=0\\
\mathbf u\times\mathbf v&\text{vector}&
\mathbf u\parallel\mathbf v\iff\mathbf u\times\mathbf v=\mathbf0
\end{array}
\]
Moreover, \(\mathbf u\times\mathbf v\) is perpendicular to \(\mathbf u\), to
\(\mathbf v\), and to every linear combination \(a\mathbf u+b\mathbf v\).
These facts are developed here as properties of vectors. They will become tools
for constructing and analysing lines and planes only in the following chapter.
\end{summarybox}